% VotaSJ --- Short Business/User White Paper (~2 pages)
% SBC LaTeX template (sbc-template.sty). See README.md for setup.
\documentclass[12pt]{article}

\usepackage{sbc-template}
\usepackage{graphicx,url}
\usepackage[utf8]{inputenc}
\usepackage[english]{babel}
\usepackage{xcolor}
\usepackage{hyperref}
\usepackage{booktabs}
\usepackage{microtype}

\hypersetup{
  colorlinks=true,
  linkcolor=black,
  citecolor=black,
  urlcolor=blue!60!black
}

% Tighten lists and section spacing to fit 2 pages.
\usepackage{enumitem}
\setlist{nosep,leftmargin=1.1em,topsep=2pt,itemsep=1pt,parsep=0pt}
\setlength{\parskip}{2pt}
\setlength{\parindent}{0pt}

\sloppy

\title{VotaSJ: Trust, Counted Live}

\author{Davi Ludvig Longen Machado\inst{1} \and
        Lucas Furlanetto Pascoali\inst{1} \\
        Arthur Clasen de Melo\inst{1} \and
        Pedro Henrique Tesman Mansani da Silva\inst{1}}

\address{Departamento de Inform\'atica e Estat\'{\i}stica (INE) \\
  Universidade Federal de Santa Catarina (UFSC) \\
  Florian\'opolis --- SC --- Brasil
  \email{\{davi.ludvig,lucas.furlanetto,arthur.clasen,pedro.tesman\}@grad.ufsc.br}}

\begin{document}
\maketitle

\begin{abstract}
\noindent VotaSJ is a blockchain-anchored participatory-budgeting platform
for S\~ao Jos\'e/SC. Citizens authenticate with gov.br, vote from their
phone, and audit the tally and the budget execution in real time --- without
having to trust the City Hall to count its own vote. The platform replaces
a BRL 300--500k physical process with a BRL 90--160k digital one, governed
by a 3-of-5 multisig (City Hall, UFSC, TCE-SC) and a 14-day timelock on
rule changes. This short paper makes the case to the institutional buyer
(\S\ref{sec:investor}) and to the citizen (\S\ref{sec:citizen}).
\end{abstract}

%% Resumo intentionally omitted: VotaSJ artifacts are English-only.

\section*{For the Investor and the Municipality}\label{sec:investor}

\noindent\textbf{Problem.} Brazilian municipalities run participatory
budgets that almost no one trusts. In S\~ao Jos\'e/SC: turnout below
\textbf{2\,\%} of 180{,}000 voters (\(<\)3{,}600 people decide tens of
millions of reais); operational cost of \textbf{BRL 300--500k/cycle}
dominated by physical assemblies and \(\approx\)3-week manual tallying;
and a tally run, counted, and announced by the City Hall itself, so every
contested result ends in an unfalsifiable accusation of manipulation.

\noindent\textbf{Solution.} A platform on Polygon PoS where citizens
authenticate with \texttt{gov.br}, vote from their phone, and audit
results live. The blockchain holds only what the City Hall must not be
allowed to rewrite --- the eligibility commitment, the votes, the tally,
and the budget-execution milestones. Personal data, proposal text, photos,
and UX stay off-chain. Administration is a \textbf{3-of-5 multisig}
(3 City Hall, 1 UFSC, 1 TCE-SC) and any contract change is gated by a
\textbf{14-day timelock}, so the rules cannot move inside an active vote
and the chain keeps running even if the City Hall stops cooperating.

\noindent\textbf{Why now.} \texttt{gov.br} is mature (\(\approx\)150\,M
citizens, free for municipal integration); Polygon per-vote gas is under
USD 0.005; OpenZeppelin + Hardhat make audited deployments affordable;
TCE-SC and MP press for auditable processes --- we arrive with the audit
built in.

\noindent\textbf{Market.}
\begin{center}\small
\begin{tabular}{l l r}
\toprule
\textbf{Segment} & \textbf{Scope} & \textbf{3y addressable} \\
\midrule
S\~ao Jos\'e/SC (anchor) & \(\approx\)180k voters & BRL 220k/yr \\
Greater Florian\'opolis & 4 cities, \(\approx\)1.1M voters & BRL 0.8--1.2M/yr \\
Santa Catarina state & \(\approx\)5.5M voters & BRL 4--6M/yr \\
BR cities \(>\)50k inhab. & \(\approx\)340 cities, \(\approx\)85M voters & BRL 60M+/yr \\
\bottomrule
\end{tabular}
\end{center}

\noindent TAM: \(>\)340 cities legally required to run participatory
cycles. SAM: \(\approx\)80 with budgeted digital
modernization. SOM (36\,mo): single-digit municipalities anchored on
S\~ao Jos\'e.

\noindent\textbf{B2G SaaS model.} Annual subscription BRL 180--250k/city;
per-cycle setup BRL 40k; implementation consulting BRL 60--120k; white-label
for professional councils. Citizens never pay. \emph{Unit economics:}
physical process \textbf{BRL 300--500k/cycle}; VotaSJ \textbf{BRL
90--160k/cycle, all-in}; net savings \textbf{BRL 170--300k/cycle}, before
counting the political value of an unimpeachable result.

\noindent\textbf{Traction.} Production-ready platform ---
\texttt{VoterRegistry}, \texttt{ParticipatoryBudget}, and
\texttt{ExecutionTracker} deployed under a UUPS proxy with a 3-of-5
multisig and 14-day timelock, exercised by a 40-test Hardhat suite,
CI-enforced lint/coverage gates, \(\approx\)430k\,gas measured for a full
citizen flow. Institutional partners: UFSC Blockchain Lab, S\~ao Jos\'e
Planning Secretariat, TCE-SC. Rollout: field pilot 2026-Q3, full S\~ao
Jos\'e cycle 2026-Q4, multi-city expansion 2027.

\noindent\textbf{Ask.} An institutional pilot agreement with S\~ao Jos\'e/SC
and a seed envelope of \textbf{BRL 400--600k} for the first external
smart-contract audit, gov.br integration certification, and the field-pilot
cycle. First paid contract closes 12\,months after disbursement; payback is
one cycle.

\section*{For the Citizen}\label{sec:citizen}

\noindent\textbf{Your vote, your eyes.} Today, when you participate in S\~ao
Jos\'e's participatory budget --- if you participate --- you go to a
night-time assembly, raise your hand, and trust that the City Hall counts
correctly and executes what won. VotaSJ changes one thing: \textbf{you can
stop trusting and start checking}.

\noindent\textbf{What changes for you.}
\begin{itemize}
  \item \textbf{Vote from your phone, on your own schedule.} No more
        assemblies you can't attend because of work, kids, or distance.
  \item \textbf{See your vote land.} The moment you cast it, the public
        tally updates. Refresh the page and watch it.
  \item \textbf{Anyone audits the count.} Journalists, neighborhood
        associations, TCE-SC, MP --- same data, same chain, same time.
  \item \textbf{Follow execution.} You're notified when your winning
        proposal moves to contract, work start, milestones, and finish,
        with evidence you can check.
  \item \textbf{Your data stays private.} CPF, name, address: none of it
        goes on the blockchain. Only a one-way hash that proves
        you're eligible, not who you are.
\end{itemize}

\noindent\textbf{Three steps.} (1) Sign in with \texttt{gov.br}
--- the federal ID you already use. (2) Pick a proposal: browse what
neighbors submitted, read the budget, vote once. (3) Watch the result and
the execution: real-time tally and per-milestone updates. No need to know
what a ``blockchain'' is. No crypto, no fee, no extra app.

\noindent\textbf{Forquilhinhas, next cycle.} Maria, a cashier whose shift
ends at 22:30, opens VotaSJ on the bus home and votes for the renovated
daycare. Months later she is notified --- contract signed, work started,
daycare reopens --- each step with verifiable photographic evidence.
90\,seconds of participation; the institution actually delivered.

\medskip
\noindent\textbf{Trust, but verify --- and now you actually can.}

\medskip
{\footnotesize\noindent\textbf{Companion document and references.} A
16-page technical white paper covers architecture, threat model, gas
measurements, team, and roadmap. Sources: \emph{Estatuto da Cidade}
(\url{https://www.planalto.gov.br/ccivil_03/leis/leis_2001/l10257.htm}),
\emph{LGPD}
(\url{https://www.planalto.gov.br/ccivil_03/_ato2015-2018/2018/lei/l13709.htm}),
\emph{gov.br} (\url{https://acesso.gov.br/roteiro-tecnico/}),
\emph{Polygon PoS} (\url{https://docs.polygon.technology/}).}

\end{document}
